\subsection{基环的扩张}

设 g 是 K 上的李代数，T 是其张量代数，J 是由 \(x \otimes y - y \otimes x - [x, y]\)（x, y 属于 g）生成的 T 的双边理想，U = T/J。设 \(K_{1}\) 是含单位元的交换环，\(\sigma\) 是将 1 映到 1 的从 K 到 \(K_{1}\) 的同态。则 \(\mathfrak{g}_{(\mathbf{K}_{1})}\) 的张量代数自然地与 \(\mathrm{T}_{(\mathbf{K}_{1})}\) 相同。令 J' 为由下式生成的 \(J'\)\(\mathrm{T}_{(\mathbf{K}_{1})}\) 的双边理想：

\[
x ^ {\prime} \otimes y ^ {\prime} - y ^ {\prime} \otimes x ^ {\prime} - [ x ^ {\prime}, y ^ {\prime} ]
\]

\((x', y' \text{ 属于 } g_{(K_1)})\)。显然，\(J_{(K_1)}\) 在 \(T_{(K_1)}\) 中的典范像包含于 J'。要看出两者相等，只需证明：若 \(J'\)\(J'\)\(x'\) 和 \(y'\) 表示 \(g_{(K_1)}\) 中的两个元素，则 \(x' \otimes y' - y' \otimes x' - [x', y']\) 属于此像。现在

\[
x ^ {\prime} = \sum_ {i} x _ {i} \otimes \lambda_ {i}, y ^ {\prime} = \sum_ {j} y _ {j} \otimes \mu_ {j} (x _ {i}, y _ {j} \text {   in   } \mathfrak {g}, \lambda_ {i}, \mu_ {j} \text {   in   } \mathrm{K} _ {1});
\]

因此

\[
x ^ {\prime} \otimes y ^ {\prime} - y ^ {\prime} \otimes x ^ {\prime} - [ x ^ {\prime}, y ^ {\prime} ] = \sum_ {i, j} (x _ {i} \otimes y _ {j} - y _ {j} \otimes x _ {i} - [ x _ {i}, y _ {j} ]) \otimes \lambda_ {i} \mu_ {j}
\]

这就证明了断言。于是可见，\(\mathrm{U}_{(\mathbf{K}_{1})} = (\mathrm{T}/\mathrm{J})_{(\mathbf{K}_{1})}\) 自然地与 \(\mathrm{T}_{(\mathbf{K}_{1})}/\mathrm{J}'\) 相同：\(\mathfrak{g}_{(\mathbf{K}_{1})}\) 的包络代数自然地与 \(\mathrm{U}_{(\mathbf{K}_{1})}\) 相同，而 \(\mathfrak{g}_{(\mathbf{K}_{1})}\) 到其包络代数的典范映射与 \(\sigma \otimes 1\) 相同（其中 \(\sigma\) 表示 g 到 U 的典范映射）。

